% R581: EGA I, 0, 2.1.1--2.1.4.
% Traducción: OpenAI Codex, GPT-6.1 Sol, esfuerzo Ultra.
\oldpage[0\textsubscript{I}]{21}

\section{Espacios irreducibles. Espacios noetherianos}
\label{section:0.2-fr}

\subsection{Espacios irreducibles}
\label{subsection:0.2.1-fr}

\begin{env}[2.1.1]
\label{0.2.1.1-fr}
Se dice que un espacio topológico $X$ es \emph{irreducible} si no es vacío
y no es unión de dos subespacios cerrados distintos de $X$. Esto
equivale a decir que $X\neq\varnothing$ y que la intersección de dos
abiertos no vacíos de $X$ (y, por tanto, de un número finito de ellos) no es
vacía, o que todo abierto no vacío es denso en todo el espacio, o que toda parte
cerrada $\neq X$ tiene \emph{interior vacío}, o, finalmente, que todo abierto de $X$ es
\emph{conexo}.
\end{env}

\begin{env}[2.1.2]
\label{0.2.1.2-fr}
Para que un subespacio $Y$ de un espacio topológico $X$ sea irreducible, es
necesario y suficiente que su adherencia $\overline Y$ sea irreducible. En
particular, todo subespacio que es la adherencia $\overline{\{x\}}$ de un
subespacio reducido a un punto es irreducible; expresaremos la relación
$y\in\overline{\{x\}}$ (equivalente a
$\overline{\{y\}}\subset\overline{\{x\}}$) diciendo que $y$ es
\emph{una especialización de $x$}, o que $x$ es \emph{una generalización de $y$}.
Cuando en un espacio irreducible $X$ existe un punto $x$ tal que
$X=\overline{\{x\}}$, diremos que $x$ es un \emph{punto genérico} de $X$.
Todo abierto no vacío de $X$ contiene entonces a $x$, y todo subespacio que contiene
a $x$ admite a $x$ como punto genérico.
\end{env}

\begin{env}[2.1.3]
\label{0.2.1.3-fr}
Recordemos que se llama \emph{espacio de Kolmogoroff} a un espacio topológico $X$
que satisface el axioma de separación siguiente:

\smallskip
\noindent
\textup{(T$_0$)} Si $x\neq y$ son dos puntos cualesquiera de $X$, existe
un conjunto abierto que contiene a uno de los puntos $x$, $y$ y no al otro.

Si un espacio de Kolmogoroff irreducible admite un punto genérico, admite
\emph{uno solo}, pues todo abierto no vacío contiene todo punto
genérico.

Recordemos que un espacio topológico $X$ se llama \emph{cuasicompacto} si de todo
recubrimiento abierto de $X$ se puede extraer un subrecubrimiento finito de $X$ (o,
lo que es equivalente, si toda familia filtrante decreciente de conjuntos
cerrados no vacíos tiene una intersección no vacía). Si $X$ es un espacio
cuasicompacto, toda parte cerrada no vacía $A$ de $X$ contiene un conjunto
cerrado no vacío \emph{minimal} $M$, pues el conjunto de las partes cerradas no vacías
de $A$ es inductivo para la relación $\supset$; si, además, $X$ es un espacio
de Kolmogoroff, $M$ se reduce necesariamente a un solo punto (o, como se
dice por abuso de lenguaje, es un \emph{punto cerrado}).
\end{env}

\begin{env}[2.1.4]
\label{0.2.1.4-fr}
En un espacio irreducible $X$, todo subespacio abierto no vacío $U$ es
irreducible, y si $X$ admite un punto genérico $x$, $x$ es también un punto
genérico de $U$.

Sea $(U_\alpha)$ un recubrimiento (con conjunto de índices no vacío) de un
espacio topológico $X$, formado por abiertos no vacíos; para que $X$ sea
irreducible, es necesario y suficiente que $U_\alpha$ sea irreducible para todo
$\alpha$, y que $U_\alpha\cap U_\beta\neq\varnothing$ para cualesquiera
$\alpha$, $\beta$. La condición es evidentemente necesaria; para ver que
es suficiente, basta probar que, si $V$ es un abierto no vacío de $X$,
$V\cap U_\alpha$ no es vacío para todo $\alpha$, pues entonces $V\cap U_\alpha$
es denso en $U_\alpha$ para todo $\alpha$, y, por tanto, $V$ es denso en
$X$. Ahora bien, existe al menos un índice $\gamma$ tal que
$V\cap U_\gamma\neq\varnothing$; por tanto, $V\cap U_\gamma$ es denso en
$U_\gamma$, y como, para todo $\alpha$,
$U_\alpha\cap U_\gamma\neq\varnothing$, también se tiene
$V\cap U_\alpha\cap U_\gamma\neq\varnothing$.
\end{env}
